\section{Cel, założenia i plan dowodu}
Słowo \emph{paradoks} oznacza tutaj twierdzenie sprzeczne z intuicją, a nie sprzeczność logiczną. Pracujemy w zwykłej teorii mnogości ZFC, czyli z aksjomatem wyboru. Kawałkiem nazywamy dowolny podzbiór kuli: nie musi być bryłą, nie musi być spójny ani mierzalny w sensie Lebesgue'a.

Zakładamy znajomość zbiorów, funkcji, liczb rzeczywistych oraz elementarnego rachunku macierzy i wektorów. Przypominamy potrzebne pojęcia grupowe; nie zakładamy znajomości grup wolnych, teorii działań grup ani teorii rozkładów paradoksalnych. Nie konstruujemy liczb rzeczywistych ani nie wyprowadzamy podstaw algebry liniowej z aksjomatów. Teoria miary pojawia się dopiero w komentarzu po dowodzie i nie jest narzędziem samej konstrukcji.

\subsection{Co dokładnie zostanie udowodnione?}
Przez $\|x\|$ oznaczamy długość euklidesową wektora $x\in\R^3$. Niech
\[
 B^3=\{x\in\R^3:\|x\|\le1\},\qquad S^2=\{x\in\R^3:\|x\|=1\}.
\]
To odpowiednio kula domknięta i jej brzeg, czyli sfera. Kula zawiera również punkty wnętrza; nie wolno jej utożsamiać ze sferą.

\begin{twierdzenie}[Banach--Tarski: podwojenie kuli]\label{tw:banach-tarski}
Istnieje skończona rodzina parami rozłącznych zbiorów $P_1,\ldots,P_N\subset B^3$ oraz ruchy sztywne $f_i(x)=Q_ix+v_i$, gdzie $Q_i\in\mathrm{SO}(3)$, takie że
\[
 B^3=\bigsqcup_{i=1}^{N}P_i,\qquad
 \bigsqcup_{i=1}^{N}f_i(P_i)
 =(B^3+(-3,0,0))\sqcup(B^3+(3,0,0)).
\]
\end{twierdzenie}
W obu zapisach unie są rozłączne. Każdy kawałek jest użyty dokładnie raz; nie rozciągamy go, nie pomniejszamy i nie tworzymy jego dodatkowej kopii. Konstrukcja podana niżej używa co najwyżej $64$ kawałków; nie jest to twierdzenie o najmniejszej możliwej liczbie kawałków.

Grupę $\mathrm{SO}(3)$, czyli grupę obrotów zachowujących orientację i punkt $0$, zdefiniujemy macierzowo w podsekcji~\ref{subsec:so3}. Dodanie $v_i$ pozwala przesuwać kawałki. Odległość między środkami dwóch kul docelowych wynosi $6>2$, więc rzeczywiście są one rozłączne.

\subsection{Siedem etapów}
\begin{enumerate}
\item Zbudujemy grupę wolną $F_2$ z liter $a,a^{-1},b,b^{-1}$.
\item Podzielimy $F_2$ na cztery zbiory, z których dwa odtworzą pierwszą kopię $F_2$, a pozostałe dwa drugą.
\item Pokażemy dwa konkretne obroty, które nie spełniają żadnej nietrywialnej relacji. Dzięki temu otrzymamy monomorfizm $F_2\to\mathrm{SO}(3)$.
\item Usuniemy ze sfery przeliczalny zbiór punktów stałych $D$. Na pozostałym zbiorze $X$ przeniesiemy rozkład grupy na każdą orbitę.
\item Wykażemy, że dodanie $D$ nie przeszkadza: sfera i $X$ są równoważne przez rozkład na dwa kawałki.
\item Przedłużymy rozkład sfery wzdłuż promieni, otrzymując rozkład kuli bez środka, a następnie osobno wchłoniemy środek.
\item Umieścimy dwie formalne kopie w dwóch różnych miejscach przestrzeni. Sprawdzimy skończoność, rozłączność i pokrycie.
\end{enumerate}
Najtrudniejszy krok algebraiczny, czyli dowód, że wskazane obroty są wolnymi generatorami, przeprowadzimy rachunkiem macierzowym modulo $3$. Nie odwołujemy się do twierdzenia o istnieniu „odpowiednich obrotów”.

